\begin{lemma}
\label{lemma-generic-finite-presentation}
Let $R \subset S$ be an inclusion of domains.
Assume that $R \to S$ is of finite type.
There exists a nonzero $f \in R$, and a nonzero $g \in S$
such that $R_f \to S_{fg}$ is of finite presentation.
\end{lemma}

\begin{proof}
We use induction on the number of generators of the original
$R$-algebra $S$. For no generators, $S=R$ and we take $f=g=1$.
For one generator write $S=R[x]/\mathfrak q$, where
$\mathfrak q\cap R=(0)$. If $\mathfrak q=(0)$, again take $f=g=1$.
Otherwise choose a nonzero polynomial $h\in\mathfrak q$ of least
degree and retain its full expression
$$
h=f x^d+a_{d-1}x^{d-1}+\cdots+a_0,
\qquad f\ne0,\quad f,a_0,\ldots,a_{d-1}\in R.
$$
Here $d\geq1$ since $R\to S$ is injective. Over $R_f$ its leading
coefficient $f$ is a unit, but the polynomial $h$ is unchanged.
Division by $h$, cancelling a leading term $b x^e$ by subtracting
$bf^{-1}x^{e-d}h$, gives a unique remainder of degree less than $d$.
Uniqueness holds because the leading coefficient of a nonzero
multiple of $h$ is a nonzero coefficient times the unit $f$.
The canonical map
$$
R_f[x]/(h)\longrightarrow S_f
$$
is surjective. If a remainder $v\in R_f[x]$ of degree less than $d$
maps to zero, choose $a\geq0$ with $f^av\in R[x]$.
Its image in $S_f$ is zero, so $f^b f^av$ maps to zero in $S$ for
some $b\geq0$. Since $S$ is a domain and $f\ne0$ in $S$, already
$f^av$ maps to zero in $S$. If $v\ne0$, the domain property of $R$
implies $f^av\ne0$ of degree less than $d$, contrary to the choice
of $h$. Thus $v=0$, proving injectivity and finite presentation.
This proves the one-generator case with the original $f$ and $g=1$.

\medskip\noindent
For $n>1$ let $x_1,\ldots,x_n$ be the original generators and
$S'=R[x_1,\ldots,x_{n-1}]\subset S$. By the induction hypothesis
there are nonzero $f\in R$ and $g\in S'$ for which
$R_f\to S'_{fg}$ is finitely presented. The extension
$S'_{fg}\subset S_{fg}$ is an inclusion of domains generated by
$x_n$. Apply the one-generator case to obtain nonzero
$f'\in S'_{fg}$ and $g'\in S_{fg}$ with
$$
(S'_{fg})_{f'}\longrightarrow (S_{fg})_{f'g'}
$$
finitely presented. Retain both elements and write their actual
fractions as $f'=a/(fg)^r$ and $g'=b/(fg)^s$, with
$a\in S'$, $b\in S$, $r,s\geq0$ and $a,b\ne0$.
Set $G=gab\in S$, which is nonzero. There are canonical inverse
$S$-algebra maps
$$
(S_{fg})_{f'g'}\ \longleftrightarrow\ S_{fG}.
$$
Indeed, in $S_{fG}$ the product $fgab$ is a unit, so each of
$f,g,a,b$ is a unit. Thus $fg,f',g'$ are units and the map from
the left ring exists. Conversely, in the left ring $fg$ and $f'g'$
are units, so $f,g,f',g'$ are units, and then
$a=f'(fg)^r$, $b=g'(fg)^s$ are units. Hence $fG=fgab$ is a unit
and the reverse map exists. The two composites fix $S$ and every
inverse adjoined in the indicated localizations, so they are identities.

\medskip\noindent
Finally $S'_{fg}\to(S'_{fg})_{f'}$ has the finite presentation
$S'_{fg}[z]/(f'z-1)$. Finite presentations compose: if
$P=A[X_1,\ldots,X_l]/(F_1,\ldots,F_v)$ and
$Q=P[Y_1,\ldots,Y_t]/(H_1,\ldots,H_w)$, lift each of the finitely
many coefficients in the $H_j$ to $A[X_1,\ldots,X_l]$ to obtain
$\widetilde H_j$. The canonical map identifies
$Q$ with $A[X_1,\ldots,X_l,Y_1,\ldots,Y_t]/(F_i,\widetilde H_j)$:
maps in both directions send every variable to its displayed class,
and their composites fix the coefficient ring and all variables.
Applying this twice proves that
$R_f\to(S_{fg})_{f'g'}\cong S_{fG}$ is finitely presented, with
the required elements $f\in R$ and $G\in S$.
\end{proof}